% Einheit 4 von 4 – Nullstellen und Schnittpunkte berechnen (bank/quadratische-funktionen/e4.jsonl, zusammenbau v0.1)
%% TODO Zweigzeile Teil 1: was hier gelernt wird (Satz in Schülersprache aus dem Einheitstitel) – Einheit 4
\einheitenkopf[e4]{Einheit 4 von 4 · Nullstellen und Schnittpunkte berechnen}
\zweigzeile{ab Kl. 10 (am Gymnasium ab Kl. 9) · P10 oft}
\setcounter{aufgabe}{16}

%% TODO Titel: Ich-kann-Satz fehlt in der Bank (Platzhalter: Kettenname) – Nr. 17
%% TODO Anweisung: ein Satz über der ersten Teilaufgabe fehlt in der Bank – Nr. 17
\begin{aufgabe}{Nullstellen und Schnittpunkte}
\begin{teile}
\teil Parabel $f(x) = x^2 - 2x$ und Gerade $g(x) = x + 4$ – welche Gleichung setzt du an? Nur hinschreiben. \leerfeld = \leerfeld
\end{teile}
%% TODO Darstellung für schwach fehlt in der Bank (quadratische-funktionen-e4-k1-s1-v1; Abschnitt „Für schwache Schüler“)
\swz{$\text{$f(x) = (x - 1)^2 - 4$ – Nullstellen?}$}{}
%% TODO Darstellung für schwach fehlt in der Bank (quadratische-funktionen-e4-k1-s1-v2; Abschnitt „Für schwache Schüler“)
\swz{$\text{$f(x) = (x + 2)^2 - 9$ – Nullstellen?}$}{}
%% TODO Darstellung für schwach fehlt in der Bank (quadratische-funktionen-e4-k1-s1-v3; Abschnitt „Für schwache Schüler“)
\swz{$\text{$f(x) = (x - 4)^2 - 1$ – Nullstellen?}$}{}
%% TODO Darstellung für schwach fehlt in der Bank (quadratische-funktionen-e4-k1-s1-v4; Abschnitt „Für schwache Schüler“)
\swz{$\text{$f(x) = (x + 3)^2 - 25$ – Nullstellen?}$}{}
%% TODO Darstellung für schwach fehlt in der Bank (quadratische-funktionen-e4-k1-s1-v5; Abschnitt „Für schwache Schüler“)
\swz{$\text{$f(x) = (x - 5)^2 - 36$ – Nullstellen?}$}{}
\swfrage{Was bleibt gleich, was ändert sich?}
%% TODO Darstellung für schwach fehlt in der Bank (quadratische-funktionen-e4-k1-s2-v1; Abschnitt „Für schwache Schüler“)
\swz[3]{$f(x) = (x - 2)^2 + 3$ – wie viele Nullstellen? Begründe ohne Rechnung.}{}
%% TODO Darstellung für schwach fehlt in der Bank (quadratische-funktionen-e4-k1-s3-v1; Abschnitt „Für schwache Schüler“)
\swz{$\text{$f(x) = x^2 - 2x - 8$ – Nullstellen?}$}{}
%% TODO Darstellung für schwach fehlt in der Bank (quadratische-funktionen-e4-k1-s4-v1; Abschnitt „Für schwache Schüler“)
\swz{$\text{$f(x) = x^2 - 4x + 1$ – Nullstellen? Runde auf zwei Stellen.}$}{}
%% TODO Darstellung für schwach fehlt in der Bank (quadratische-funktionen-e4-k1-s5-v1; Abschnitt „Für schwache Schüler“)
\swz{$\text{$f(x) = 3x^2 - 3x - 18$ – Nullstellen?}$}{}
%% TODO Darstellung für schwach fehlt in der Bank (quadratische-funktionen-e4-k1-s6-v1; Abschnitt „Für schwache Schüler“)
\swz{$\text{$f(x) = x^2 + x + 1$ – für welche $x$ ist $f(x) = 7$?}$}{}
%% TODO Darstellung für schwach fehlt in der Bank (quadratische-funktionen-e4-k1-s7-v1; Abschnitt „Für schwache Schüler“)
\swz{$\text{$f(x) = x^2 + x - 4$ und $g(x) = 2x + 2$ – an welchen Stellen $x$ schneiden sich Parabel und Gerade?}$}{}
%% TODO Darstellung für schwach fehlt in der Bank (quadratische-funktionen-e4-k1-s8-v1; Abschnitt „Für schwache Schüler“)
\swz{$\text{$f(x) = (x - 2)^2 - 1$ und $g(x) = x - 3$ – an welchen Stellen $x$ schneiden sich Parabel und Gerade?}$}{}
%% TODO Darstellung für schwach fehlt in der Bank (quadratische-funktionen-e4-k1-s9-v1; Abschnitt „Für schwache Schüler“)
\swz{$\text{$f(x) = x^2 - 4x + 1$ und $g(x) = 2x - 4$ – Schnittpunkte?}$}{}
%% TODO Darstellung für schwach fehlt in der Bank (quadratische-funktionen-e4-k1-s10-v1; Abschnitt „Für schwache Schüler“)
\swz{$f(x) = x^2 - 5$, $g(x) = 2x + 3$ – ist $P(4|11)$ ein Schnittpunkt?}{}
%% TODO Darstellung für schwach fehlt in der Bank (quadratische-funktionen-e4-k1-s11-v1; Abschnitt „Für schwache Schüler“)
\swz{$f(x) = (x - 2)^2 + 1$ – gib eine Gerade $g$ an, die mit der Parabel keinen Punkt gemeinsam hat.}{}
%% TODO Darstellung für schwach fehlt in der Bank (quadratische-funktionen-e4-k1-s12-v1; Abschnitt „Für schwache Schüler“)
\swz{$\text{$f(x) = x^2 + 3x - 2$ und $h(x) = x^2 - x + 6$ – Schnittpunkt?}$}{}
\swz[3]{$\text{Die Parabel $f(x) = (x - 4)^2 - 6$ und die Gerade $g(x) = -2x + 5$ schneiden sich in zwei Punkten. Berechne ihre Koordinaten. \hfill (P10 2022 OS)}$}{}
\swz[3]{$\text{Berechne die Schnittpunkte der Parabel $f(x) = x^2 - 10x + 22$ mit der $x$-Achse. Runde auf zwei Stellen. \hfill (P10 2025 OS)}$}{}
\swz[3]{$\text{Die Parabel $f(x) = -x^2 + 4x + 2$ ist nach unten geöffnet. Für welche $x$ ist $f(x) = -10$? \hfill (P10 2023 OS)}$}{}
\end{aufgabe}

%% TODO Titel: Ich-kann-Satz fehlt in der Bank (Platzhalter: Kettenname) – Nr. 18
%% TODO Anweisung: ein Satz über der ersten Teilaufgabe fehlt in der Bank – Nr. 18
\begin{aufgabe}{Nullstellen und Schnittpunkte – Fehler finden, Begründen, Darstellungswechsel, Anwendung}
\swz[3]{$\text{Lina berechnet die Nullstellen von $f(x) = (x + 1)^2 - 25$: \rechnung{0 &= (x + 1)^2 - 25 \\ 25 &= (x + 1)^2 \\ x + 1 &= 5 \\ x &= 4} Finde den Fehler und korrigiere.}$}{}
\swfrage{Warum liegt jede Nullstelle einer Parabel auf der $x$-Achse?}
\swz{$f(x) = (x - 2)^2 - 3$ und $g(x) = x - 3$ – lies die Schnittpunkte am Graphen ab und bestätige einen durch Einsetzen.}{\begin{ksys}[xmin=-1,xmax=5,ymin=-4,ymax=2,ablesen]
\parabel{1}{2}{-3}{f}
\gerade{1}{-3}{g}
\end{ksys}}
\swz[3]{$\text{Ein Ball fliegt auf der Bahn $h(x) = -0{,}1(x - 6)^2 + 4{,}9$ ($x$: Entfernung in m, $h$: Höhe in m). Wie weit fliegt der Ball, bis er auf dem Boden aufkommt?}$}{\begin{ksys}[xmin=0,xmax=14,ymin=0,ymax=6,xstep=2,xlabel=$x$ in m,ylabel=$h$ in m]
\funktionab{-0.1*(\x-6)^2+4.9}{h}{0}{13}
\end{ksys}}
\end{aufgabe}

\uebersichtskasten{Nullstellen: 0 = p(x) setzen – die Null tritt an die Stelle von p(x), der Term bleibt, wo er stand. Aus der Scheitelpunktform durch Wurzelziehen, aus der Normalform mit der p-q-Formel (→ Quadratische Gleichungen). \\ 0 = (x − 2)² − 16 → 16 = (x − 2)² → x − 2 = 4 oder x − 2 = −4 → x₁ = 6, x₂ = −2 \quad Nullstellen N₁(6 | 0), N₂(−2 | 0) \\ Schnittpunkte von Gerade und Parabel: Terme gleichsetzen p(x) = f(x), alles auf eine Seite bringen, quadratische Gleichung lösen, y-Werte mit der Geraden berechnen. \\ x² − 6x + 10 = x + 4 → x² − 7x + 6 = 0 → x₁ = 1, x₂ = 6 \quad y über f: S₁(1 | 5), S₂(6 | 10) \\ Zwei Lösungen: zwei Schnittpunkte. Eine Lösung: Berührpunkt. Keine Lösung (Wurzel aus einer negativen Zahl): kein gemeinsamer Punkt.}
